% Ерөнхий дүгнэлт

\phantomsection
\addchaptertocentry{Ерөнхий дүгнэлт}
\label{Summary}
%-------------------------------------------------------------------------------
%	Summary
%-------------------------------------------------------------------------------
\section*{Ерөнхий дүгнэлт}

\begin{sloppypar}
\paragraph{}
``Ахлах сургуулийн сурагчдад зориулсан сургуулийн зөвлөгчийн веб систем хөгжүүлэх'' сэдэвтэй дипломын ажлын зорилго нь ахлах сургуулийн сурагчдын гадаадын их, дээд сургуульд элсэх бэлтгэл үйл явцыг нэгдсэн платформоор удирдах, зөвлөх багш болон сурагчдын хоорондын харилцаа, даалгаврын менежмент, баримт бичгийн солилцоог нэгтгэсэн вебд суурилсан системийг хөгжүүлэхэд оршиж байв. Тус системийг хөгжүүлснээр зөвлөх багшийн ажлын ачааллыг бууруулах, сурагчдын гүйцэтгэлийн хяналтыг сайжруулах, элсэлтийн үйл явцыг илүү ил тод, зохион байгуулалттай болгох нөхцөл бүрдсэн.

\par Төслийн хөгжүүлэлтийн явцад хэрэглэгчийн болон системийн шаардлагуудыг тодорхойлж, UML диаграммуудыг гаргасан. Системийн архитектур нь хэрэглэгчтэй харилцах хэсэг болох front-end болон системийн логик үйл ажиллагааг хангах хэсэг болох back-end гэсэн 2 хэсгээс бүрдэж, хоорондоо REST API ашиглан холбогдох Modular Monolith архитектурын зарчмаар зохиомжлогдсон. Front-end хэсгийг Next.js 15 болон React дээр суурилан хөгжүүлж, Tailwind CSS ашиглан responsive дизайн хэрэгжүүлсэн. Back-end хэсэгт Next.js API Routes ашиглан Auth, User, Task, File, Chat, Notification, AI Recommendation, Payment гэсэн 8 модулийг бие даасан байдлаар зохион байгуулсан. 
Өгөгдлийг MongoDB Atlas (AWS Seoul бүс) дээр хадгалж, хэрэглэгчийн баталгаажуулалтыг JWT token болон httpOnly cookie ашиглан аюулгүй хэрэгжүүлсэн.

\par Системийн гол боломжуудын нэг болох AI сургуулийн зөвлөмжийн модулийг OpenAI GPT-4o-mini загвар ашиглан хэрэгжүүлсэн. Сурагчийн GPA, IELTS, SAT болон гадуурх үйл ажиллагааны мэдээлэлд үндэслэн тохирох сургуулиудыг ``Safety'', ``Match'', ``Reach'' ангиллаар автоматаар зөвлөдөг энэхүү модуль нь системийн хамгийн цоо шинэ хувь нэмэр болж чадсан. Мөн сурагчдын essay, transcript, CV, recommendation letter зэрэг баримт бичгүүдийг Vercel Blob үйлчилгээнд хадгалах, хувилбарын түүхийг хянах, чат харилцааг polling архитектурын тусламжтайгаар дэмжих зэрэг функциональ боломжуудыг бүрэн хэрэгжүүлсэн. Даалгаврын хугацааны сануулгыг Vercel Cron болон Nodemailer ашиглан өдөр бүр автоматаар и-мэйлээр илгээх боломж нэмэгдсэн.

\par \textbf{Төслийн менежмент болон кодын удирдлага.}
Хөгжүүлэлтийн явцад Git version control систем ашиглан кодын өөрчлөлтийг удирдсан. Эх код нь GitHub репозиторт байршуулагдсан бөгөөд feature branch стратеги ашиглан бие даасан модулиудыг тусдаа салбарт хөгжүүлж, code review хийсний дараа үндсэн салбарт нэгтгэсэн. GitHub репозиторыг Vercel платформтай шууд холбосон тул \texttt{main} салбарт өөрчлөлт оруулах бүрт автомат CI/CD дамжлагаар шинэ хувилбар байршуулагдах боломж бий болсон. Энэ нь хөгжүүлэлтийн цикл хурдасч, нэвтрүүлэлтийн алдаа эрс буурахад хувь нэмэр оруулсан.

\par \textbf{Системийн байршуулалт.}
Хөгжүүлэгдсэн систем нь Vercel платформ дээр байршуулагдан ажиллаж байна. Vercel нь hosting, reverse proxy болон serverless function орчинг нэгтгэн санал болгодог тул тусдаа сервер удирдах шаардлагагүйгээр системийн бүх бүрэлдэхүүн хэсгийг нэгдсэн байдлаар байршуулах боломж олгосон. Өгөгдлийн сан болох MongoDB Atlas-ийг AWS Seoul бүсэд байрлуулснаар Монголоос хандах сүлжээний хоцролтыг бууруулсан. Файл хадгалалтыг Vercel Blob-ийн тусламжтайгаар өгөгдлийн сангаас тусгаарлан зохион байгуулсан. Системд хандахын тулд Vercel-ийн өгч буй домайн нэрийг ашиглаж байгаа бөгөөд HTTPS холболтыг Vercel автоматаар хангаж байна. Цаашид тусгай домайн нэр авч, бүрэн нийтийн хэрэглээнд зориулсан байдлаар байршуулахаар төлөвлөж байна.

\par \textbf{Бодит хэрэглэгчид нэвтрүүлэлт.}
Систем нь хөгжүүлэлтийн эцсийн шатанд бодит нөхцөлд туршигдсан. Бета тестийн явцад зөвлөх багш болон сурагчид оролцуулан тестийн үйлдлүүдийг гүйцэтгүүлж, дараагийн хэсэгт дурдсан үр дүнг олж авсан. Одоогийн байдлаар систем нь прод орчинд байрлуулагдан ажиллаж байгаа бөгөөд анхны хэрэглэгчдийн бүлгийг онбординг хийх шатанд байна. Цаашид хэрэглэгчдийн тоог нэмэгдүүлж, Монгол дахь гадаадад суралцах хүсэлтэй сурагчдыг дэмжих зөвлөх байгууллагуудтай хамтран ажиллахаар төлөвлөж байна.

\par Энэхүү дипломын ажил нь ахлах сургуулийн сурагчдын гадаадын их, дээд сургуульд элсэх бэлтгэлийн үйл явцыг нэгдсэн платформоор удирдах боломжийг бүрдүүлсэн системийн бичиг баримт болон бүрэн ажиллагаатай веб системийг багтаасан болно. Цаашдаа системийг сурагчдын тоо болон сургуулиудын хамрах хүрээг өргөтгөх, мобайл апп хэлбэрт шилжүүлэх, AI зөвлөмжийн нарийвчлалыг дээшлүүлэх зэрэг чиглэлээр хөгжүүлэх боломжтой.
\end{sloppypar}
